\documentclass[a4paper]{article}
% Espanol (Mexico) con XeLaTeX: fontspec para fuentes Unicode, polyglossia
% para la division silabica y los nombres del documento en la variante mexicana.
\usepackage{fontspec}
\usepackage{polyglossia}
\setdefaultlanguage[variant=mexican]{spanish}
% polyglossia no traduce los nombres de listings.
\AtBeginDocument{\providecommand{\lstlistingname}{}\renewcommand{\lstlistingname}{Listado}%
  \providecommand{\lstlistlistingname}{}\renewcommand{\lstlistlistingname}{Índice de listados}}
\usepackage{amsmath,amssymb}
\usepackage{graphicx}
\usepackage[margin=2.35cm]{geometry}
\usepackage[most]{tcolorbox}
\usepackage{etoolbox}
\usepackage{listings}
\usepackage{xcolor}
\usepackage{hyperref}
\usepackage{booktabs,longtable,tabularx,array}
\usepackage{subcaption}
\usepackage{float}
\usepackage{enumitem}
\usepackage{microtype}
\usepackage{fancyhdr}
\usepackage{needspace}

\definecolor{kbBack}{HTML}{F2F6FA}
\definecolor{kbFrame}{HTML}{9DB4CC}
\definecolor{kbTitle}{HTML}{52708F}
\definecolor{ibBack}{HTML}{FBF5E7}
\definecolor{ibFrame}{HTML}{C9A961}
\definecolor{ibTitle}{HTML}{7D5F1E}
\definecolor{wbBack}{HTML}{FCF2F0}
\definecolor{wbFrame}{HTML}{D6A096}
\definecolor{wbTitle}{HTML}{9E4F43}
\definecolor{dbBack}{HTML}{F4F7F3}
\definecolor{dbFrame}{HTML}{AEC2AC}
\definecolor{dbTitle}{HTML}{5F7F64}

\tcbset{softbox/.style={
  enhanced,breakable,boxrule=0.8pt,arc=2pt,left=3mm,right=3mm,
  top=2mm,bottom=2mm,coltitle=white,fonttitle=\bfseries,
  toptitle=0.6mm,bottomtitle=0.6mm
}}
\newtcolorbox{knowledgebox}[1]{softbox,colback=kbBack,colframe=kbFrame,colbacktitle=kbTitle,title={#1}}
\newtcolorbox{importantbox}[1]{softbox,colback=ibBack,colframe=ibFrame,colbacktitle=ibTitle,title={#1}}
\newtcolorbox{warningbox}[1]{softbox,colback=wbBack,colframe=wbFrame,colbacktitle=wbTitle,title={#1}}
\newtcolorbox{dialoguebox}[1]{softbox,colback=dbBack,colframe=dbFrame,colbacktitle=dbTitle,title={#1}}

\lstset{
  language=Python,basicstyle=\ttfamily\small,
  keywordstyle=\color{kbTitle}\bfseries,stringstyle=\color{wbTitle},
  commentstyle=\color{dbTitle},breaklines=true,frame=single,numbers=left,
  numberstyle=\tiny\color{gray},captionpos=b,columns=fullflexible,
  keepspaces=true,showstringspaces=false,extendedchars=true
}
\BeforeBeginEnvironment{lstlisting}{\Needspace{12\baselineskip}}

\hypersetup{colorlinks=true,linkcolor=kbTitle,urlcolor=kbTitle,citecolor=wbTitle}
\setlist{itemsep=0.25em,topsep=0.4em}
\setlength{\parindent}{2em}
\setlength{\parskip}{0.25em}
\newcolumntype{L}[1]{>{\raggedright\arraybackslash}p{#1}}

\providecommand{\notetitle}{Notas del curso: [tema del curso]}
\providecommand{\noteauthors}{Elaboradas a partir de los materiales públicos de Stanford CS153}
\providecommand{\notedate}{Primavera de 2026}
\providecommand{\videochannel}{Stanford Online}
\providecommand{\videopublishdate}{2026}
\providecommand{\videoduration}{[duración del video]}
\providecommand{\videourl}{}
\providecommand{\videocoverpath}{cover.jpg}
\providecommand{\coursesite}{https://cs153.stanford.edu/}
\providecommand{\repourl}{https://github.com/hqhq1025/ai-course-notes}
\providecommand{\skillversion}{wdkns/wdkns-skills@39f1a04}

\newcommand{\figsource}[1]{\par\vspace{-0.3em}{\footnotesize\color{gray}\emph{Fuente: #1}}\par}
\newcommand{\teachervoice}[1]{\begin{knowledgebox}{Nota de clase}#1\end{knowledgebox}}
\newcommand{\term}[1]{\textbf{#1}}

\pagestyle{fancy}
\fancyhf{}
\fancyfoot[C]{\thepage}
\fancyfoot[R]{\footnotesize\color{gray}\href{\repourl}{AI Course Notes}}
\renewcommand{\headrulewidth}{0pt}
\renewcommand{\footrulewidth}{0pt}

\newcommand{\makecscover}{
\begin{titlepage}
\centering
\vspace*{0.7cm}
{\Large Stanford CS153 · Frontier Systems\par}
\vspace{0.8cm}
{\huge\bfseries \notetitle\par}
\vspace{0.6cm}
{\large \noteauthors\par}
\vspace{0.25cm}
{\large \notedate\par}
\vspace{0.9cm}
\ifdefempty{\videocoverpath}{}{
  \includegraphics[width=0.86\textwidth,height=0.43\textheight,keepaspectratio]{\videocoverpath}\par
}
\vfill
\begin{tcolorbox}[width=0.92\textwidth,colback=black!2!white,colframe=black!45,boxrule=0.7pt,arc=2pt]
\textbf{Autor o canal del video}: \videochannel\par
\textbf{Fecha de publicación}: \videopublishdate\par
\textbf{Duración del video}: \videoduration\par
\textbf{Sitio del curso}: \href{\coursesite}{\nolinkurl{\coursesite}}\par
\ifdefempty{\videourl}{}{\textbf{Enlace del video}: \href{\videourl}{\nolinkurl{\videourl}}\par}
\textbf{Normas de elaboración}: \skillversion\ + las normas source-first / teacher-voice / visual-QA de este repositorio
\end{tcolorbox}
\end{titlepage}
\tableofcontents
\newpage
}


\renewcommand{\notetitle}{Lecture 08: Sistemas de seguridad infantil en línea: detección, revisión, reporte y Safety by Design}
\renewcommand{\noteauthors}{Elaborado a partir de la entrevista con Julie Cordua, los subtítulos con marcas de tiempo y los materiales oficiales de Thorn/NCMEC/NIST}
\renewcommand{\notedate}{Curso de invierno de 2025 · versión reescrita source-first de 2026}
\renewcommand{\videochannel}{Publicación de cursos anteriores de Stanford CS153}
\renewcommand{\videopublishdate}{2025-02-12}
\renewcommand{\videoduration}{37:53}
\renewcommand{\videourl}{https://www.youtube.com/watch?v=MBD0Ah9cpYU}
\renewcommand{\videocoverpath}{cover.jpg}

\newcommand{\lecturefigure}[3]{%
\begin{figure}[H]
  \centering
  \includegraphics[width=0.96\textwidth]{images/#1}
  \caption{#2}
  \figsource{#3}
\end{figure}
}

\begin{document}
\makecscover

\section{Auditoría de fuentes: convertir un tema sensible en una clase de sistemas verificable}

Esta sesión la imparte Julie Cordua, CEO de Thorn, y trata de cómo la tecnología de seguridad infantil en línea ayuda a las plataformas a identificar riesgos, reducir la difusión repetida, organizar la revisión humana y entregar las pistas que cumplen los requisitos al sistema legal de reporte e investigación. Al consultarlo el 11 de agosto de 2026, el video público pedía iniciar sesión; el repositorio conserva la portada oficial y 833 subtítulos con marcas de tiempo, por lo que este texto toma los subtítulos de la clase como hilo narrativo y contrasta los mecanismos permanentes con los materiales oficiales de Thorn, NCMEC y NIST, sin presentar capacidades de producto posteriores como hechos de la clase de 2025.

\begin{warningbox}{Tratamiento del contenido sensible y límites de la evidencia}
Este texto no relata, describe ni muestra contenido ilegal, ni ofrece métodos para generarlo, difundirlo o evadir su detección. Los casos individuales mencionados en clase se abstraen únicamente como eventos del sistema, por ejemplo: «la plataforma detectó un archivo de alto riesgo y completó el reporte». El número de reportes, el número de archivos, las proporciones de los estudios y la escala de los productos dependen del año, de la definición y del método de deduplicación; salvo que se indique de forma explícita, las cifras de años distintos no deben restarse ni compararse directamente.
\end{warningbox}

\begin{importantbox}{Los tres ciclos cerrados de esta sesión}
\begin{enumerate}
  \item \textbf{Response loop}: carga o mensaje $\rightarrow$ detección $\rightarrow$ prioridad $\rightarrow$ revisión humana $\rightarrow$ medidas de la plataforma y reporte legal;
  \item \textbf{Learning loop}: resultados verificados y autorizados $\rightarrow$ trusted data/hash/model update $\rightarrow$ cobertura de riesgos futuros;
  \item \textbf{Governance loop}: nuevas tecnologías y nuevos abusos $\rightarrow$ threat research $\rightarrow$ Safety by Design $\rightarrow$ pruebas, auditoría y transparencia.
\end{enumerate}
\end{importantbox}

\begin{knowledgebox}{Términos clave: no mezclar cuatro tipos de objeto bajo «contenido malo»}
\textbf{Content} es el archivo o mensaje que procesa la plataforma; \textbf{signal} es una hash match, un score del modelo o un patrón de comportamiento; \textbf{case} es un evento revisable organizado por cuenta, tiempo y contexto; \textbf{outcome} es el resultado en forma de medidas de la plataforma, reporte legal, investigación o protección de la víctima. Toda métrica de ingeniería debe indicar qué capa mide; de lo contrario, «aumentó el volumen de detección» puede significar una mejor cobertura, pero también solo archivos repetidos, un cambio de umbral o un cambio en el criterio de reporte.
\end{knowledgebox}

\subsection{Por qué esto no es «entrenar un clasificador»}

Esta sección define primero el objeto de ingeniería. \term{CSAM} es la sigla de child sexual abuse material y designa el material de abuso sexual infantil tal como lo definen la ley y las políticas de las plataformas; el término describe a la víctima y la evidencia, y no debe confundirse con «contenido para adultos» común. Un modelo solo puede producir una match o un score; el resultado real de protección depende además de la policy de la plataforma, la cola de moderación, la preservación de la evidencia, los organismos legales de reporte, la investigación policial y la identificación de las víctimas. Si falla cualquier eslabón, la precisión del eslabón anterior no se convierte automáticamente en un resultado de seguridad.

\lecturefigure{01-response-pipeline.png}{La tecnología de seguridad infantil es una cadena de respuesta de extremo a extremo formada por detect, review, action/report y victim identification.}{Subtítulos locales 00:00--05:50; rediseño conceptual no sensible basado en el flujo presentado en clase.}

\begin{knowledgebox}{Lectura de la figura: distinguir primero entre signal, decision y outcome}
Lo que producen el classifier o el hash service de la izquierda es una \textbf{señal}; la moderación de la plataforma y la policy dan lugar a una \textbf{decisión de producto}; los organismos de reporte y los investigadores buscan un \textbf{resultado de investigación} dentro del procedimiento legal. El score de un modelo no puede probar por sí solo una identidad, una intención subjetiva ni una responsabilidad legal. El diseño del sistema debe registrar quién ejecuta cada transición de estado, con base en qué policy, qué audit evidence deja y cómo se apelan o corrigen los errores.
\end{knowledgebox}

\teachervoice{Cordua no abre la sesión con casos impactantes para captar la atención, sino con una pregunta: cuando la carga de reportes de las plataformas no deja de crecer, ¿puede la tecnología conectar de verdad la detección, las medidas adoptadas y el tiempo limitado de los investigadores? Esta pregunta lleva la clase de la «identificación de contenido malo» a las end-to-end operations.}

\subsection{Resumen de la sección}

El objetivo de esta sesión es explicar el sistema de respuesta completo, no hacer promoción de producto para un modelo en particular. Todos los mecanismos que siguen deben responder tres preguntas: qué tipo de riesgo reducen, a quién le transfieren el trabajo y dónde están los límites de su evidencia.
\section{La escala no es un número, sino varias colas que no pueden mezclarse}

La sección anterior estableció la cadena de respuesta; esta sección explica por qué el «crecimiento de la escala» desborda a las organizaciones. Las estadísticas públicas suelen presentar al mismo tiempo reports, files, accounts, incidents y victims, cuyas unidades, tasas de duplicación y criterios anuales difieren. Si un equipo de ingeniería escribe en el dashboard solo la cifra más grande, estimará mal la carga de almacenamiento, revisión, investigación y entrenamiento de modelos.

\subsection{Reports, files, duplicates y novel candidates}

Un report puede contener varios archivos, y un mismo archivo puede ser reportado varias veces por distintas plataformas o cuentas; lo que los investigators realmente necesitan atender con prioridad suele ser la parte que puede corresponder a una nueva víctima, a un riesgo persistente o a material antes desconocido. La \term{revictimization} (nuevo daño) se refiere al daño repetido que sufre la víctima cuando el material se copia, se ve o se difunde de forma continua; por ello, incluso la reaparición de un «archivo conocido» merece bloquearse con rapidez, pero no debe confundirse con el «número de casos nuevos».

\lecturefigure{02-scale-funnel.png}{El volumen de reportes, el volumen de archivos, el contenido duplicado, los candidatos desconocidos y la escasa capacidad de revisión son medidas distintas.}{Subtítulos locales 00:00--05:50; nota sobre los criterios anuales de NCMEC; redibujo conceptual.}

\begin{knowledgebox}{Lectura de la figura: primero pregunte por el denominador y las reglas de deduplicación}
Al leer cualquier dato de escala, primero escriba la unidad: un report es un envío, un file es un adjunto, un candidate es un objeto preseleccionado por un algoritmo y pendiente de revisión, y un case es la unidad de investigación que una institución forma al agregar elementos según sus reglas. También hay que preguntar si se deduplica por cryptographic hash, perceptual hash, account, incidente o victim. NCMEC ha ajustado sus métodos de estadística y de reporte en distintos años, por lo que las páginas de 2023, 2024 y 2025 solo pueden interpretarse dentro de sus propias definiciones, y un punto de inflexión en la curva no puede atribuirse directamente a un cambio interanual del daño real.
\end{knowledgebox}

\begin{importantbox}{La planeación de capacidad debe considerar cuatro tasas a la vez}
Sea $\lambda_r$ el número de reportes que llegan por unidad de tiempo, $m_f$ el número promedio de archivos por reporte, $p_k$ la proporción de contenido conocido y $p_u$ la proporción de candidatos desconocidos que requieren revisión humana; entonces la tasa aproximada de llegada a revisión es
\[
\lambda_h \approx \lambda_r m_f p_u.
\]
Aquí $\lambda_h$ no es el número de víctimas ni el número de casos legales; es solo una estimación de la carga de la cola de revisión. Las coincidencias con contenido conocido siguen activando acciones de atención y reportes, por lo que el hecho de que desaparezcan de la cola humana no significa que no tengan costo.
\end{importantbox}

\teachervoice{En clase se distinguió repetidamente entre la «detección y reporte» de las plataformas y la «identificación de víctimas y persecución de los responsables» de los investigadores. El ponente enfatizó que lo que menos necesitan las instituciones al final de la cadena es otro montón de archivos sin priorizar, sin contexto y con gran cantidad de duplicados.}

\subsection{Resumen de la sección}

En esencia, el problema de escala es de colas heterogéneas: la difusión duplicada debe bloquearse a gran velocidad, el contenido desconocido debe revisarse con cuidado, las pistas de alta prioridad deben llevar contexto y los investigadores necesitan un case package accionable. Una sola cifra total no puede orientar el diseño del sistema.
\section{Defensa ante contenido conocido: capacidades y límites del Hash Matching}

Cuando una institución confiable ya verificó un archivo y este se incorporó conforme a la ley al reference set, la defensa escalable más madura es el matching. Evita que el modelo o una persona tengan que volver a interpretar el mismo contenido cada vez y es adecuado para atender la propagación repetida a gran escala; sin embargo, solo puede identificar contenido «lo bastante idéntico a un objeto conocido», no puede descubrir todo el contenido nuevo y tampoco puede sustituir la evaluación de legalidad ni el proceso de reporte de la plataforma.

\subsection{Cryptographic hash y perceptual hash}

\term{cryptographic hash} (hash criptográfico) asigna a una secuencia de bytes un resumen de longitud fija y es adecuado para determinar si un archivo es idéntico byte por byte; una compresión o un recorte leves cambian el resumen. \term{perceptual hash} (hash perceptual) extrae la estructura visual o auditiva, de modo que medios similares que pasaron por transformaciones comunes todavía pueden quedar en representaciones cercanas; mejora la recuperación de variantes, pero requiere umbrales, reference provenance y control de las coincidencias erróneas.

\lecturefigure{03-hash-matching.png}{El Hash matching es la defensa mínima para bloquear a gran escala material conocido y verificado.}{Subtítulos locales 05:10--06:00; mecanismo oficial de Thorn Safer Match; redibujo conceptual.}

\begin{knowledgebox}{Lectura de la figura: Match no es «volver a identificar el contenido»}
La plataforma primero calcula el hash o el fingerprint del incoming media y luego lo compara con un reference set controlado; si hay coincidencia, el caso pasa a la ruta de policy, revisión o reporte de la plataforma. Un cryptographic match suele aportar evidencia más sólida de byte identity, mientras que un perceptual match expresa similitud. En ambos casos hay que registrar hash algorithm/version, reference source, threshold, match confidence y las acciones humanas posteriores, para poder reproducir y auditar los resultados después de actualizar el sistema.
\end{knowledgebox}

\begin{warningbox}{Una Hash database no es una blocklist común}
Los Reference hash están asociados a evidencia de sensibilidad extremadamente alta, por lo que deben restringirse las fuentes de importación, los permisos de acceso, la capacidad de exportación y la visibilidad de los registros. La plataforma no puede ignorar la membership inference, el abuso interno, el etiquetado erróneo, la filtración entre tenants ni las obligaciones de eliminación por el hecho de que «solo guarda hashes». El servicio de matching debe devolver el resultado mínimo necesario, en lugar de exponer el contenido de las reference o una interfaz enumerable.
\end{warningbox}

\subsection{Por qué la coincidencia con contenido conocido aún requiere intervención humana y políticas}

Esta sección reconecta la señal de hash con la cadena de respuesta. Un mismo archivo puede aparecer en la obtención de evidencia para una denuncia, en reportajes periodísticos, en investigación, en sistemas legítimos de las autoridades o en el entorno de carga de un usuario, y en cada caso el contexto y las obligaciones legales pueden ser distintos; las distintas jurisdicciones, tipos de servicio y funciones de producto también tienen requisitos de reporte diferentes. Por ello, el equipo de ingeniería debe modelar el «match» como una señal de alta confianza, y no como una orden universal de eliminación que se salte toda policy.

\begin{importantbox}{Campos mínimos de un evento de coincidencia auditable}
Un durable match event contiene como mínimo: content identifier, hash algorithm/version, reference list/version, match distance, tenant/account, event time, policy version, review/action, reporting handoff y retry state. Los campos sensibles deben almacenarse en capas según el principio de mínimo privilegio, y deben diseñarse procesos explícitos para la eliminación, el legal hold y la incident response.
\end{importantbox}

\subsection{Resumen de la sección}

El Hash matching obtiene velocidad y consistencia a cambio de depender de un reference set verificado, y es la base de la known-material defense. Sus límites son igual de claros: no puede cubrir contenido desconocido, no puede inferir la identidad ni la intención de quien actúa, ni exime de la revisión del contexto ni de los procedimientos legales.
\section{Defensa ante contenido desconocido: Predictive Detection, calibración y revisión humana}

Tanto las herramientas generativas como el contenido recién subido producen objetos que el reference set no ha visto; por eso, la siguiente capa es la predictive detection. Lo esencial aquí no es «si el modelo puede asignar una puntuación», sino cómo convertir el score en una cola de revisión operable y sostenible: el umbral debe ajustarse según el riesgo y la capacidad, la salida debe pasar por calibration y la decisión final debe conservar la human review y los límites de apelación.

\subsection{Del classifier score al verified outcome}

Un \term{predictive classifier} (clasificador predictivo) aprende, a partir de los datos de entrenamiento, la relación estadística entre la entrada y las etiquetas de riesgo, y produce una clase o un score. La \term{human review} (revisión humana) consiste en que personal capacitado, sujeto a restricciones de permisos y a medidas de protección de su salud, tome decisiones con base en el contexto y la policy; no se trata de simplemente «echar un vistazo al resultado del modelo». La salida del modelo es una triage signal, no una conclusión jurídica.

\lecturefigure{04-predictive-classifier.png}{La Predictive AI amplía la cobertura más allá de los hash conocidos, pero debe pasar por la cola y la revisión.}{Subtítulos locales 10:10--13:40; mecanismo oficial de Thorn Safer Predict; redibujo conceptual.}

\begin{knowledgebox}{Lectura de la figura: el significado del score depende del entrenamiento y la calibración}
El modelo produce primero un raw score; después, la plataforma asigna cada objeto a un nivel de prioridad según el threshold; el reviewer ve el case context filtrado por permisos y escribe el verified outcome de vuelta en el sistema controlado. La magnitud del raw score no puede compararse directamente entre model versions, tipos de medio o entornos de cliente. Antes del despliegue debe precisarse si el score es un ranking, una estimated probability o un uncalibrated margin, y deben registrarse model/version y threshold/version.
\end{knowledgebox}

\teachervoice{Cordua usa la evolución de Safer Match a Safer Predict para explicar que, cuando los atacantes o los modelos generativos crean objetos nuevos sin cesar, ya no basta con buscar solo hash conocidos; pero el nombre «Predict» tampoco significa que el sistema pueda prescindir de las personas. El punto central de la clase es ampliar la cobertura de candidatos y adelantar los elementos que más vale la pena revisar.}

\subsection{Precision, recall, calibration y queue capacity}

Esta sección reúne en una misma figura las métricas del modelo y las restricciones operativas. La \term{precision} (precisión) indica qué proporción de los objetos que el sistema marca como positivos se verifica como positiva; el \term{recall} (exhaustividad) indica qué proporción de todos los positivos reales encuentra el sistema; la \term{calibration} (calibración) indica si la probabilidad predicha coincide con la frecuencia real observada a largo plazo. Ninguna de las tres puede optimizarse por separado sin considerar la base rate, el label process y la reviewer capacity.

\lecturefigure{10-evaluation-queue.png}{La calidad del modelo y la capacidad de revisión deben diseñarse juntas; el umbral es una decisión operativa, no una constante pura del modelo.}{Subtítulos locales 10:10--16:20; enfoque de medición de riesgos del NIST AI RMF; redibujo conceptual.}

\begin{importantbox}{Definición de las métricas y explicación de los símbolos}
Si se denotan los verdaderos positivos, falsos positivos y falsos negativos como $TP,FP,FN$, respectivamente, entonces
\[
\mathrm{precision}=\frac{TP}{TP+FP},\qquad
\mathrm{recall}=\frac{TP}{TP+FN}.
\]
La precision determina cuántos elementos inválidos encuentra el revisor entre las alertas, y el recall determina cuántos riesgos reales omite el sistema. Si el número diario de candidatos es $N$, la tasa de alertas con el umbral dado es $q(\tau)$, la capacidad diaria de cada revisor es $c$ y el número de revisores es $R$, una operación estable requiere al menos
\[
Nq(\tau) \lesssim Rc,
\]
donde $\tau$ es el umbral. Esta desigualdad es una restricción de capacidad; no significa que el objetivo de seguridad se reduzca a «evitar que la cola se acumule».
\end{importantbox}

\begin{warningbox}{Una accuracy alta puede ocultar una base rate extremadamente baja}
En un entorno donde los positivos reales son escasísimos, un modelo que casi siempre predice negativo puede tener una accuracy muy alta. La plataforma debe reportar precision/recall estratificados por tipo de medio, región, idioma, escenarios relacionados con la edad y punto de entrada del producto, y monitorear la distribution shift. Un cambio de umbral afecta los falsos negativos, pero también altera la exposición de los revisores y la carga de los reportes obligatorios por ley.
\end{warningbox}

\subsection{La moderator triage es parte del sistema de seguridad}

La revisión humana no es un último paso ilimitado ni gratuito. El \term{triage} (triaje/priorización) ordena los candidatos según gravedad, confianza, urgencia y posibilidad de actuar; las buenas herramientas, además, ocultan la información visual no esencial, agrupan los objetos duplicados, ofrecen el contexto mínimo necesario y dan soporte a la rotación de turnos y a mecanismos de salud mental. Al leer la figura de esta sección hay que observar a la vez la «calidad de la decisión» y la «exposición humana»; no basta con perseguir el rendimiento.

\lecturefigure{05-moderator-triage.png}{El triage debe maximizar el valor de protección y, al mismo tiempo, minimizar la exposición sensible innecesaria de los revisores.}{Subtítulos locales 13:40--16:20; redibujo conceptual basado en la discusión de Safer Review.}

\begin{knowledgebox}{Lectura de la figura: un case package es mejor que una cola de archivos sueltos}
La capa de entrada aplica primero hash, ordenamiento por modelo y agrupación de duplicados; el workbench muestra solo la evidencia mínima necesaria para aplicar la policy; la reviewer decision y su rationale se registran de forma estructurada; después, el case package incorpora la account, el tiempo, los eventos relacionados y el action history. Así se reducen las visualizaciones repetidas y, además, la segunda revisión, el muestreo de calidad, los reportes y la incident review pueden realizarse sobre un mismo objeto auditable.
\end{knowledgebox}

\begin{warningbox}{El reviewer wellbeing no puede depender de la «fortaleza psicológica»}
La organización necesita exposure minimization, capacitación, rotación, descansos, apoyo profesional, monitoreo de accesos y clear escalation. Los falsos positivos del modelo tienen un costo humano real; por eso, un false positive no es solo una cifra en el dashboard. A la inversa, tampoco se puede ajustar el threshold hasta provocar omisiones sistemáticas con tal de reducir la exposición. Las políticas de producto, de modelo y de personal deben asumir el riesgo en conjunto.
\end{warningbox}

\subsection{Resumen de la sección}

El criterio de éxito de la predictive detection no es la accuracy en un leaderboard, sino que una calibrated signal ayude a los revisores a atender antes los objetos de alto riesgo dentro de una cola sostenible. Las métricas del modelo, el umbral, el case packaging, la salud del personal y los mecanismos de apelación deben validarse en conjunto.
\section{Trusted Data: gobernanza de datos para modelos de alto riesgo}

La sección anterior dependía de etiquetas de entrenamiento y de verified outcome; esta sección indaga de dónde provienen esos datos, quién puede acceder a ellos y cómo demostrar que son legales y confiables. Los datos de seguridad de alto riesgo son a la vez escasos y sensibles, por lo que no pueden seguir el método de la extracción común de internet ni el de «primero copiar al data lake y después gobernar». La barrera central de la capacidad de datos no es la cantidad de archivos, sino la autorización, la provenance, la segregation, el label protocol y la auditoría continua.

\subsection{Provenance, segregation y retention}

La \term{provenance} (origen y linaje de procesamiento) registra quién aportó los datos, bajo qué autorización y mediante qué transformaciones llegaron a una versión de entrenamiento o de evaluación. La \term{data segregation} (aislamiento de datos) separa los datos de sensibilidad extrema de los registros generales, los entornos de desarrollo, las cuentas de bajos privilegios y otros inquilinos. La \term{retention} (política de conservación) establece por qué se conservan los datos, durante cuánto tiempo y cuándo se eliminan o pasan a legal hold, en lugar de conservarlos de manera permanente por omisión.

\lecturefigure{06-trusted-data.png}{Los datos de entrenamiento de alto riesgo requieren fuentes confiables, procesamiento trazable, acceso aislado y una política de conservación explícita.}{Subtítulos locales 16:20--20:30; contexto de la colaboración entre Thorn y NCMEC; redibujo conceptual.}

\begin{knowledgebox}{Lectura de la figura: la gobernanza de datos abarca todo el lifecycle}
La trusted source demuestra la autorización y el uso; la provenance vincula la evidencia original, las transformaciones de desidentificación, la label y la dataset version; la segregation restringe entornos, redes y personal; la retention/audit se ocupa de la eliminación, la respuesta a incidentes y las inspecciones regulatorias. Si falta cualquiera de estas piezas, el equipo podría no ser capaz de explicar por qué un modelo aprendió cierto patrón, por qué un empleado pudo acceder a los datos, o si una solicitud de eliminación realmente alcanzó a los derived artifacts.
\end{knowledgebox}

\begin{importantbox}{Label quality no equivale a votación por mayoría}
Las etiquetas de alto riesgo requieren límites claros fijados por definiciones legales, la policy de la plataforma y las guías operativas. El conjunto de entrenamiento debe distinguir entre verified positive, policy violation, uncertain, benign hard negative y unavailable context; la evaluación además debe registrar el reviewer agreement y la adjudication. Mezclar todo el contenido «incómodo» en una sola clase positiva produce modelos imposibles de interpretar y amplía los daños colaterales contra la expresión legítima.
\end{importantbox}

\teachervoice{El ponente enfatiza que Thorn puede entrenar este tipo de modelos no porque haya encontrado un conjunto de datos público más grande, sino porque, mediante una colaboración confiable de largo plazo, obtuvo material controlado, legal y con contexto. Nota de clase: en dominios sensibles, el data access es en sí mismo una responsabilidad de gobernanza, no una ventaja común de adquisición.}

\subsection{Los datos de entrenamiento, evaluación y producción deben tener funciones separadas}

Esta sección separa los usos de los datos para evitar que el equipo del modelo ajuste parámetros una y otra vez sobre el mismo reference set y aun así presente el resultado como una prueba independiente. Los datos de entrenamiento sirven para ajustar los parámetros; los datos de validation, para elegir umbrales y modelos; los datos de test, para la evaluación final; y el production feedback solo puede pasar a la siguiente versión después de analizar sus sesgos y sus mecanismos de selección. Para cada conjunto deben fijarse el snapshot, la versión, los roles de acceso y las verificaciones de contamination.

\begin{warningbox}{El verified feedback también puede tener selection bias}
Solo los objetos que el modelo envía a revisión tienden a recibir etiquetas, por lo que los datos de retroalimentación se inclinan de manera natural hacia las regiones que el modelo actual considera «sospechosas». Si se reentrena directamente con esas muestras, el sistema podría volverse cada vez más ciego a los riesgos que quedan fuera de los umbrales anteriores. El equipo necesita muestreo aleatorio, shadow evaluation, un reference set independiente y red-team evidence externa para descubrir los blind spots.
\end{warningbox}

\subsection{Resumen de la sección}

Los trusted data son la combinación de autorización legal, aislamiento técnico, protocolos de etiquetado y evidencia de auditoría. Si un modelo de alto riesgo no puede explicar el origen, el uso y la ruta de eliminación de sus datos de entrenamiento, no cumple las condiciones para un despliegue seguro, aunque sus métricas fuera de línea sean muy altas.
\section{Generative AI: adelantar el control de riesgos a todo el ciclo de vida}

Lo que cambia la IA generativa no es solo que «se produzcan archivos nuevos»: también cambian la suplantación de identidad, las variantes de contenido, la integración en aplicaciones y la velocidad de distribución. Escanear únicamente en el punto final de carga a la plataforma traslada toda la responsabilidad al extremo de la cadena de suministro. Un enfoque más sólido es distribuir el threat model entre training data, model access, application, hosting/distribution e incident response, y verificarlo de forma continua en las tres etapas de development, deployment y maintenance.

\subsection{La threat surface abarca cuatro capas}

Esta sección parte de una perspectiva de cadena de suministro. Los desarrolladores del modelo controlan el entrenamiento y el fine-tuning; el servicio del modelo controla el access y el abuse monitoring; los desarrolladores de aplicaciones controlan el prompt flow, la user identity y los guardrail del producto; las plataformas de distribución controlan el hosting, el sharing, el reporting y la removal. Cada participante ve señales distintas, por lo que no se puede exigir que la plataforma del extremo recupere por sí sola la provenance o las relaciones de cuenta que ya se perdieron aguas arriba.

\lecturefigure{07-genai-threat-surface.png}{Los riesgos de la Generative AI abarcan las capas de desarrollo, acceso al modelo, aplicación y distribución.}{Subtítulos locales 05:45--10:20, 20:30--24:10; libro blanco Safety by Design de Thorn; redibujo conceptual.}

\begin{knowledgebox}{Lectura de la figura: escribir la capability y la obligation de cada capa}
La capa de development puede realizar dataset curation, evaluation y release gating; la capa de model access puede realizar authentication, rate limit, usage monitoring y abuse response; la capa de application puede diseñar una UX age-appropriate, contextual guardrail y user reporting; la capa de distribution puede ejecutar detection, removal, evidence preservation y lawful reporting. Ninguna capa posee todos los hechos, pero cada una debe conservar las señales mínimas de seguridad que puedan transmitirse a la capa siguiente.
\end{knowledgebox}

\begin{warningbox}{No convertir el threat model en un manual de ataque}
Los equipos de seguridad necesitan entender las categorías de abuso, los puntos de entrada y las señales observables, pero los apuntes públicos, los informes de pruebas y la documentación del producto no deben ofrecer pasos de evasión reproducibles, prompts sensibles ni flujos de generación. Una transparencia eficaz debe explicar la cobertura, los métodos de prueba, las categorías de falla, los tiempos de respuesta y los límites de gobernanza, en lugar de divulgar detalles que reduzcan de forma significativa el costo del abuso.
\end{warningbox}

\subsection{Safety by Design es una disciplina del ciclo de vida}

\term{Safety by Design} consiste en incorporar de forma continua el control de riesgos durante el diseño, el desarrollo, el despliegue y el mantenimiento, en lugar de añadir un moderation endpoint después del lanzamiento. Exige que el equipo escriba el threat model al definir los objetivos del producto, que realice evaluation/red teaming antes del entrenamiento y del release, que diseñe monitoring, escalation y user controls al desplegar, y que atienda el drift, los incident y el transparency reporting durante el mantenimiento.

\lecturefigure{08-safety-lifecycle.png}{Safety by Design abarca design, development, deployment y maintenance, no una única revisión de lanzamiento.}{Subtítulos locales 20:30--26:20; principios de Thorn/All Tech Is Human y NIST AI 600-1; redibujo conceptual.}

\begin{knowledgebox}{Lectura de la figura: el release gate es solo un punto intermedio}
La etapa de design define el intended use, el misuse y los grupos afectados; la etapa de development gobierna los datos, el modelo y las pruebas; la etapa de deployment configura el threshold, la human review, el appeal y el incident channel; la etapa de maintenance monitorea el drift, actualiza la policy, corrige vulnerabilidades y publica los avances. Un modelo que aprobó el launch checklist, si no cuenta con telemetry continua ni con responsables, puede fallar al cabo de unas semanas por cambios en el tráfico, el idioma o la integración con el producto.
\end{knowledgebox}

El \term{red teaming} (pruebas de equipo rojo) consiste en buscar activamente los modos de falla del sistema en un entorno autorizado y controlado; debe abarcar las salidas del modelo, los flujos de trabajo del producto, el abuso de cuentas, la cadena de suministro y el response process. Los resultados de las pruebas deben traducirse en owner, severity, fix, retest y residual risk, y no limitarse a mostrar un conjunto de capturas de pantalla de «inducción exitosa del modelo».

\teachervoice{Cordua describe los compromisos de la industria como un lenguaje común: los expertos en seguridad infantil no pueden reescribir el producto completo de cada empresa, pero sí pueden convertir los riesgos recurrentes en principles, standards y progress reporting que se puedan probar, de modo que la responsabilidad entre en el flujo normal de ingeniería.}

\subsection{Resumen de la sección}

La IA generativa requiere responsabilidad a lo largo de la cadena de suministro y control durante todo el ciclo de vida. El escaneo en el extremo sigue siendo importante, pero no sustituye la gobernanza de datos aguas arriba, el control de acceso al modelo, los guardrail de la aplicación, las pruebas continuas ni la transparencia. El valor de Safety by Design está en convertir la seguridad de un proyecto puntual en un sistema que se puede mantener.
\section{Integración con la plataforma: rápida, idempotente y con capacidad de auditoría}

Los modelos y los principios solo producen efecto cuando entran en el upload/message path real. Esta sección trata la capacidad de seguridad infantil como un platform service de alto riesgo: necesita integración por API, por lotes o por stream; necesita reintentos y degradación; necesita una tenant isolation estricta, y además debe conectar los estados de revisión, resolución y reporte en una state machine rastreable. La baja latencia es importante, pero no puede obtenerse una disponibilidad aparente a cambio de una silent failure.

\subsection{Separar el detection service de las acciones de la plataforma}

La plataforma puede activar la detección antes del upload, después del almacenamiento, antes de la distribución pública o después de un reporte de un usuario; cada punto tiene un latency budget y un context distintos. El detection service devuelve match/score e información de versión; el policy engine combina la superficie del producto, la región, el historial de la cuenta y las reglas de apelación de los usuarios para decidir hold, limit, remove, review o no action. Separar ambos facilita que las actualizaciones del modelo y los cambios de policy se auditen de forma independiente.

\lecturefigure{09-platform-integration.png}{La integración con la plataforma requiere cuatro capas de estado: detection, policy, action/report y audit.}{Subtítulos locales 03:20--05:50, 10:10--16:20; Thorn platform solutions; rediseño conceptual.}

\begin{knowledgebox}{Lectura de la figura: los reintentos deben ser idempotentes}
\term{idempotency} (idempotencia) significa que ejecutar repetidamente el mismo evento no produce efectos secundarios adicionales. El upload event debe llevar un event ID estable; el detection result se deduplica por content/model/version; el action service evita eliminaciones o notificaciones duplicadas; el reporting workflow evita que el mismo reporte legal se envíe varias veces por un network retry. Cada capa debe registrar request, result, retry y final disposition para poder reproducir lo ocurrido después de un incidente.
\end{knowledgebox}

\begin{importantbox}{Una máquina de estados operable}
El estado de un objeto puede escribirse como
\[
\texttt{received}\rightarrow\texttt{screened}\rightarrow
\texttt{queued}\rightarrow\texttt{reviewed}\rightarrow
\texttt{actioned}\rightarrow\texttt{reported/closed}.
\]
Cada arista la activa un actor y una policy version explícitos; las fallas pasan a retry, dead-letter o manual escalation, en lugar de tratar un timeout como benign. Un objeto puede tener en paralelo un content state, un account state y un report state; no se puede representar todo su ciclo de vida con un solo campo booleano.
\end{importantbox}

\subsection{Failure mode y estrategias de degradación}

Esta sección pasa de la ruta normal a las fallas: puede ocurrir que la detección de un tercero no esté disponible, un queue backlog, la reversión de un model version, una falla en un reference update, retraso en el audit log o un reporting endpoint timeout. La plataforma debe definir de antemano los límites de aplicación de fail-open, fail-closed, hold-for-review y rate-limited sampling, y permitir que el equipo on-call vea la acumulación, la antigüedad, la gravedad y la estimación de recuperación.

\begin{warningbox}{Tanto fail-open como fail-closed pueden perjudicar a los usuarios}
Fail-open puede permitir que objetos de alto riesgo sigan propagándose; fail-closed puede bloquear de forma masiva a usuarios legítimos o causar resoluciones irreversibles. Una solución más robusta suele estratificar según la confianza de la señal, la superficie del producto y el daño potencial: un known match de alta confianza puede pasar a un hold estricto, una predictive signal de baja confianza entra en una cola controlada, y ante un outage sistémico se activan la limitación de tasa, la publicación diferida y la escalación manual.
\end{warningbox}

\teachervoice{En clase se describe a la plataforma como la principal responsable de detection, removal y reporting, y no se presenta al proveedor de tecnología como una autoridad policial. Este límite exige que la integración del producto devuelva señales explicables y, al mismo tiempo, que la plataforma cliente conserve su propia policy, su relación con los usuarios y sus obligaciones legales.}

\subsection{Resumen de la sección}

La capacidad de seguridad infantil debe integrarse como un durable workflow, no como una sola llamada síncrona a una API. Solo al separar en capas detection, policy, action, reporting y audit puede el sistema mantener evidencia consistente durante reintentos, reversiones, degradaciones e investigaciones.
\section{Privacy, Encryption y señales mínimas necesarias}

Los sistemas de seguridad procesan contenido muy sensible e información de cuentas. La privacidad no es una «restricción externa» opuesta a la seguridad, sino un objetivo de arquitectura. El \term{end-to-end encryption} (cifrado de extremo a extremo, E2EE) significa que solo los extremos de la comunicación tienen las claves del contenido de los mensajes; ni el transporte ni el servidor pueden leer el texto en claro. Protege a los usuarios de que la plataforma, los atacantes y los intermediarios espíen el contenido, y también limita el content scanning del lado del servidor. Un diseño responsable debe precisar qué señales siguen disponibles, quién puede acceder a ellas y cómo evitar que la vigilancia se amplíe.

\subsection{Cuatro capas: content, metadata, endpoint y governance}

Esta sección no busca una solución universal para «eludir el cifrado», sino que analiza el problema por capas. Los \term{metadata} (metadatos) son las propiedades de la comunicación o del sistema distintas del contenido, por ejemplo la hora, la cuenta, el dispositivo, las relaciones de conexión y los eventos del servicio; pueden ser muy sensibles para la privacidad, y no deben recopilarse sin límite solo porque no forman parte del cuerpo del mensaje. Los endpoint controls se sitúan en el dispositivo del usuario o en el límite de la aplicación; las network signals describen comportamientos coordinados, y la governance restringe el acceso, el legal process, la appeal y la oversight.

\lecturefigure{11-privacy-layers.png}{La privacidad y la seguridad infantil requieren controles por capas de content, metadata, endpoint y governance.}{Subtítulos locales 28:40--31:10; principios de NIST y de gobernanza de plataformas; redibujo conceptual.}

\begin{knowledgebox}{Lectura de la figura: cada señal debe declarar su purpose limitation}
La content signal puede ser la más directa, pero también la más invasiva; los metadata pueden revelar relaciones anómalas, pero es fácil que afecten por error conductas sociales ordinarias; la endpoint intervention puede alertar antes al usuario, pero hay que evitar que se amplíe la vigilancia a nivel de dispositivo; la governance determina quién puede consultar, cuándo se elimina la información y cómo se apela. Para cada capa deben especificarse el purpose, la base legal, los campos mínimos, la retention, el access role, la audit y los secondary use prohibidos.
\end{knowledgebox}

\begin{warningbox}{«La privacidad y la seguridad son importantes» no es una respuesta de diseño}
Una respuesta de ingeniería real necesita un threat model y límites: si se altera la confidencialidad de extremo a extremo, si se introduce una capacidad de escaneo que pueda reutilizarse con otros fines, quién corrige los falsos positivos, si los menores de edad y los grupos vulnerables asumen riesgos adicionales, y si el gobierno o el personal interno pueden ampliar las consultas. Cualquier propuesta debe someterse a una evaluación de seguridad independiente, a una evaluación de impacto en los derechos y a una gobernanza transparente, en lugar de depender solo de las promesas del proveedor.
\end{warningbox}

\teachervoice{Cordua reconoce que el encryption impone limitaciones reales, y a la vez subraya que la plataforma debe seguir buscando vías de seguridad legítimas y que protejan la privacidad. Lo más valioso de la clase no es una postura política concreta, sino la idea de «no fingir que el trade-off no existe, y tampoco dejar por ello de diseñar el sistema».}

\subsection{Resumen de la sección}

La privacidad y la seguridad requieren mínimo privilegio, purpose limitation y una gobernanza auditable. El E2EE protege la confidencialidad del contenido y cambia lo que el servidor puede ver; los metadata, las endpoint signals o las network signals solo pueden usarse como señales de riesgo por capas, y no se convierten automáticamente en justificación de una vigilancia amplia.
\section{Del objeto individual a la conversación y la red: señales más tempranas pero más inciertas}

La detección de archivos ocurre por lo general después de que el contenido ya se produjo o se subió. Para intervenir antes, las plataformas analizan el conversation context, el comportamiento de las cuentas y los network pattern; estas señales pueden ayudar a descubrir riesgos de coordinación, pero también es más fácil que confundan relaciones normales con conductas dañinas. Esta sección subraya que «más temprano» significa «más incierto», por lo que se requiere más contexto, umbrales escalonados y escalamiento a personas.

\subsection{La text-risk detection necesita contexto}

\term{grooming-risk detection} (detección del riesgo de grooming) consiste en identificar riesgos potenciales a partir de la evolución de la conversación, el contexto relacionado con la edad, las pruebas de límites y los patrones de escalamiento, en lugar de condenar directamente con base en una sola palabra. El modelo de texto debe manejar el idioma, la jerga, la ironía, las conversaciones entre pares y el contenido educativo legítimo; su salida debe activar avisos, friction protectora o revisión especializada, no inferir automáticamente la identidad del usuario ni su intención delictiva.

\lecturefigure{12-text-risk.png}{La detección de riesgo en texto requiere contexto de la conversación, umbrales escalonados, escalamiento y revisión humana.}{Subtítulos locales 31:10--32:50; redibujo conceptual basado en la discusión en clase sobre el text classifier.}

\begin{knowledgebox}{Lectura de la figura: la unidad es la conversation trajectory}
La context window reúne el orden temporal, las relaciones entre participantes y los eventos del producto; el risk model produce un score que cambia con el tiempo; la escalation layer elige, según la edad, la confianza y la reversibilidad, entre un aviso, una restricción o una review; la human decision integra además la policy completa. Una palabra clave en una sola frase puede ser una joke, contenido educativo, una petición de ayuda o una cita; solo la trajectory y el contexto del producto reducen los errores de clasificación, y aun así no eliminan la incertidumbre.
\end{knowledgebox}

\begin{warningbox}{Tanto los false positive como los false negative necesitan una vía de reparación}
Un falso positivo puede hacer que un adolescente pierda apoyo legítimo, activar restricciones de cuenta innecesarias o exponer conversaciones privadas; un falso negativo puede hacer que se pierda la ventana de intervención temprana. El sistema debe preferir intervention reversibles y mínimas, ofrecer user reporting, appeal y specialist escalation, y evaluar la performance del modelo por separado según idioma, grupo de edad y escenario de producto.
\end{warningbox}

\teachervoice{El ponente describe la clasificación de texto como una herramienta para «ver el pattern antes», y aclara que no es una prevención completa. Advertencia de la clase: cuanto más se acerca el modelo a inferir intenciones y relaciones, más humildad necesita; el producto debe convertir la incertidumbre en protección escalonada, no en conclusiones definitivas más rápidas.}

\subsection{El graph analysis identifica conductas coordinadas}

\term{graph analysis} (análisis de grafos) representa cuentas, dispositivos, contenido, mensajes y eventos como nodos y aristas, y busca conductas coordinadas en la estructura de conexiones, los patrones repetidos y la sincronía temporal. Una sola cuenta puede no aportar evidencia suficiente, pero cuando varias cuentas comparten infraestructura, objetivos o secuencias de comportamiento se forma una señal más fuerte. El modelo de grafos también debe evitar confundir familias, escuelas, comunidades o redes públicas con grupos maliciosos.

\lecturefigure{13-network-analysis.png}{El riesgo de coordinación suele ser más evidente en el grafo de cuentas y en los patrones de comportamiento que en un objeto individual.}{Subtítulos locales 31:50--33:10; redibujo conceptual basado en la discusión sobre los network pattern de las plataformas.}

\begin{knowledgebox}{Lectura de la figura: las graph feature exigen prudencia causal}
Account/device/content forman un grafo heterogéneo; la pattern layer puede extraer características como shared device, rapid creation, contacto repetido o synchronized action; después, el coordination score entra en la investigation queue. Una conexión solo indica correlación; no prueba una relación de control ni una intención maliciosa. NAT, los dispositivos públicos, el uso compartido en familia y las migraciones pueden generar un alto grado de conexión, por lo que se necesita validación cruzada con múltiples señales y un evidence bundle explicable.
\end{knowledgebox}

\begin{importantbox}{Objeto, conversación y red son tres perspectivas complementarias}
La object-level signal sirve para evaluar riesgos conocidos o visuales; la conversation-level signal explica la evolución de las relaciones; la network-level signal descubre coordinación y reincidencia. Las tres deben compartir tiempo y provenance dentro de un case model unificado, pero conservar por separado sus umbrales, versiones de modelo y permisos de acceso. Mezclarlas en un único opaque risk score impide que el reviewer entienda o corrija los errores.
\end{importantbox}

\subsection{Resumen de la sección}

Las señales más tempranas provienen de la trayectoria de las conversaciones y de la estructura de la red, y también conllevan mayores riesgos de privacidad y de errores de clasificación. Las plataformas deben adoptar vías de intervención multiseñal, escalonadas y reversibles, y mantener siempre el score separado de cualquier conclusión jurídica o de identidad.
\section{De la detección a la prevención: Intervention Ladder y capacidades del usuario}

La mayoría de los sistemas anteriores identifican y atienden el riesgo después de que aparece; esta sección amplía el objetivo hacia la prevention. Prevenir no significa «predecir todo lo malo», sino diseñar la friction adecuada para cada grado de certeza: educación del usuario, privacidad predeterminada, adecuación a la edad, accesos para denunciar, avisos de riesgo, restricciones de contacto, hold de contenido y escalamiento a especialistas. Cuanto más irreversible es una acción, mayores son la evidencia y la revisión que requiere.

\subsection{La intervention ladder de cuatro niveles}

Una ladder práctica puede escalar de forma gradual de educate a disrupt, protect y, por último, investigate/report. La education ayuda a los usuarios a reconocer la manipulación y los canales para pedir ayuda; la disruption eleva el costo del contacto masivo, de los mensajes de desconocidos o de la migración rápida a otra plataforma; la protection ofrece configuraciones predeterminadas, block/report y trusted-contact support; solo los eventos de alta confianza pasan a revisión especializada, a acciones de la plataforma y a los procedimientos legales.

\lecturefigure{14-prevention-ladder.png}{La prevención requiere una combinación por niveles de educación, friction de producto, escalamiento protector y respuesta legal.}{Subtítulos locales 33:10--35:30; NCMEC NetSmartz y discusión en clase; reconstrucción conceptual.}

\begin{knowledgebox}{Lectura de la figura: cuanto más débil es la evidence, más reversible debe ser la acción}
Las señales de baja confianza justifican mostrar avisos de seguridad, limitar que desconocidos encuentren al usuario u ofrecer block/report; el riesgo medio puede añadir friction, retrasar el envío o pedir una confirmación adicional; solo los eventos de alta confianza y ya revisados pasan a acciones obligatorias y a reportes. Este orden protege a los usuarios de los falsos positivos y permite que el sistema siga ofreciendo ayuda desde las etapas tempranas. El software puede reducir las oportunidades y el tiempo de respuesta, pero no puede sustituir a la familia, la escuela, las líneas de ayuda ni los servicios profesionales.
\end{knowledgebox}

\teachervoice{Cordua resume la dirección como el paso de la reactive detection a la prevention, pero no afirma que la tecnología pueda resolver el problema por sí sola. La alfabetización digital de padres y jóvenes, los canales claros para pedir ayuda y las protecciones predeterminadas de la plataforma siguen siendo parte del sistema, y no «etapas manuales que pueden eliminarse una vez que el modelo entra en producción».}

\subsection{Resumen de la sección}

Un sistema de prevención debe hacer corresponder la solidez de la evidencia con la reversibilidad de la acción. La education, la friction, el user control, la specialist review y la lawful response forman juntos la línea de defensa; ningún classifier por sí solo puede cubrir toda esta ladder.
\section{Startup Minimum Viable Safety}

Las grandes plataformas pueden formar equipos dedicados de trust-and-safety. Las empresas emergentes, en cambio, suelen enfrentar riesgos desde que aparece el contenido generado por usuarios, pero les faltan experiencia, presupuesto y prioridad para atenderlos. Esperar a «gobernar cuando haya más escala» convierte en deuda técnica el modelo de datos, los registros, el flujo de reportes y la cultura. Esta sección presenta la minimum viable safety: no se trata de copiar la organización de una gran empresa, sino de establecer un ciclo cerrado mínimo que pueda ampliarse.

\subsection{Risk inventory, known-content control e incident readiness}

Primero, el equipo de una empresa emergente debe enumerar varios puntos: quién puede interactuar con quién dentro del producto, qué contenido puede subirse, si los menores de edad pueden verlo, cómo se propagan los mensajes privados y los grupos, y qué modelos o servicios de almacenamiento de terceros participan en el procesamiento. Después, elige la combinación mínima de coincidencia con contenido conocido, reportes de usuarios, moderation queue, account action, audit log y emergency escalation. También define de forma explícita qué se resolverá con build, con buy o con specialist partnership.

\lecturefigure{15-startup-maturity.png}{Las empresas emergentes necesitan pasar del risk inventory a una ruta de minimum viable safety que pueda ampliarse.}{Subtítulos locales 35:20--37:20; redibujo conceptual basado en la discusión de la clase sobre la startup barrier.}

\begin{knowledgebox}{Lectura de la figura: primero, que la responsabilidad y los incidentes tengan a dónde ir}
El risk inventory responde cuál es la superficie de exposición del producto. Los minimum controls aportan known-hash, report y basic review. El managed program agrega predictive triage, policy, metrics y auditoría. Por último, la proactive prevention incorpora threat research, red teaming y protección de usuarios. En la primera etapa, lo más importante no es comprar todas las herramientas. Lo que importa es designar un owner, conservar los eventos necesarios, establecer una vía de reporte externo y, ante un incident real, poder localizar a la persona de guardia y el registro de decisiones.
\end{knowledgebox}

\begin{importantbox}{Minimum viable safety checklist}
\begin{enumerate}
  \item Un safety owner con poder de decisión y una on-call escalation;
  \item el risk inventory del producto, el modelo de edad e interacción y la política de uso aceptable;
  \item una ruta de report/block/help visible y fácil de usar para los usuarios;
  \item known-content control, basic queue, case log y deduplicación de eventos repetidos;
  \item minimización de datos, permisos de acceso, retention e incident playbook;
  \item SLA de proveedores, degradación ante fallas, reportes legalmente obligatorios y lista de contactos de instituciones especializadas.
\end{enumerate}
\end{importantbox}

\begin{warningbox}{«Hacerlo después» amplifica el costo de migración}
Si desde el inicio no hay ID estables de user/account/content, más adelante no es posible agregar los eventos. Sin audit log, no se pueden explicar las medidas que se tomaron antes. Sin un age-aware product model, la única opción es imponer restricciones a la fuerza una vez terminada la UI. Sin supplier contract, es difícil rastrear los datos sensibles. Cuanto más tarde se incorpore la safety architecture, más probable es que obligue a reescribir identity, storage, messaging y el moderation core.
\end{warningbox}

\teachervoice{Quien expone resume los obstáculos de las startups en tres factores: awareness, money y priority. Con esto evita limitarse a culpar a los fundadores. Para la ingeniería, la clase concluye que hay que ofrecer infraestructura compartida, una ruta de madurez clara y servicios asequibles. Así, los equipos pueden empezar con un ciclo cerrado mínimo y no tienen que elegir entre «no hacer nada» y «construir por su cuenta un departamento de seguridad completo».}

\subsection{Resumen de la sección}

La seguridad en una empresa emergente empieza por tener claros los riesgos del producto, los responsables y las rutas de reporte e incidentes. Después se incorporan de forma gradual la detección, la evaluación y la prevención. Integrar esto desde el principio suele costar menos que limpiar después de un incidente, y también facilita conservar los límites correctos de datos y permisos.
\section{Responsabilidades en el ecosistema: plataformas, proveedores de tecnología, organismos de denuncia e investigadores}

Por último, volvemos al sistema de extremo a extremo. En el ecosistema de seguridad infantil ningún participante tiene todas las capacidades: la plataforma controla la relación con los usuarios y las action del producto; los proveedores de tecnología aportan las detection/review primitives; los organismos de denuncia previstos por la ley reciben las pistas y las complementan con más información; las autoridades y los servicios de atención a víctimas investigan y protegen conforme a la ley. Tener responsabilidades claras no es eludirlas. Sirve para evitar que el proveedor del modelo se exceda en sus atribuciones, que la plataforma subcontrate su gobernanza o que los investigadores se vean rebasados por datos de baja calidad.

\subsection{Límites de responsabilidad y ciclo de retroalimentación}

La plataforma debe hacerse cargo de la policy, las notificaciones a usuarios, las apelaciones, las medidas sobre cuentas y las obligaciones legales. Los proveedores de tecnología como Thorn deben ofrecer herramientas validadas y auditables, y proteger los datos sensibles. Los organismos de denuncia como el NCMEC procesan los reportes y las pistas según el papel que les asigna la ley. Las autoridades se encargan de la investigación, las pruebas y los procedimientos legales. Las organizaciones de apoyo a víctimas brindan protección continua. Un verified outcome solo puede retroalimentar al modelo y al reference system si hay autorización, minimización de datos y limitación de la finalidad.

\lecturefigure{16-ecosystem-roles.png}{Los resultados en seguridad infantil dependen de una división clara del trabajo entre plataformas, proveedores de tecnología, organismos de denuncia e investigadores.}{Subtítulos locales 37:10--37:53; reconstrucción conceptual basada en la closing synthesis de la clase.}

\begin{knowledgebox}{Lectura de la figura: cada handoff necesita un contract}
El contract entre platform y provider establece las entradas, las salidas, el SLA, el uso de los datos y los incident; el contract entre platform y reporting agency establece los campos que exige la ley, el manejo de duplicados y la cadena de custodia de las pruebas; el handoff de la agency a los investigators requiere prioritization y jurisdiction; cuando el investigation outcome regresa, su uso debe quedar restringido. Un handoff ambiguo genera reportes duplicados, campos faltantes, casos que nadie atiende o accesos que exceden las atribuciones.
\end{knowledgebox}

\begin{importantbox}{Implicaciones de ingeniería de un mission-driven business model}
Una misión sin fines de lucro y el software comercial no entran en conflicto por naturaleza: los productos de pago pueden financiar la investigación y el desarrollo continuos, el soporte, las auditorías y una infraestructura de alta disponibilidad, mientras que la gobernanza basada en la misión limita la elección de clientes, el uso de los datos y las prioridades del producto. El criterio real de aceptación no es la etiqueta de la organización, sino si puede mantener herramientas de seguridad a largo plazo, hacer públicos sus límites, aceptar supervisión y hacer que la métrica principal sea el resultado en protección y no el engagement.
\end{importantbox}

\teachervoice{El cierre de Cordua describe los papeles sin rodeos: la plataforma debe detectar y tomar medidas, y los investigadores deben identificar a las víctimas y dar seguimiento a los casos. La función de la tecnología es que la información que pasa de unos a otros sea más rápida, menos repetida y más útil para actuar, no que un modelo sustituya a las instituciones y a los profesionales.}

\subsection{Resumen de la sección}

La seguridad infantil es un sistema que abarca varias organizaciones. Con roles, data contracts, audit trails y límites de retroalimentación claros, cada parte puede aprovechar su especialidad. Si las responsabilidades son ambiguas, el riesgo se traslada al eslabón que tiene menos contexto o menos recursos.
\section{Síntesis y ampliación}

Esta clase amplía la seguridad infantil en línea: deja de tratarla como un problema aislado de clasificación de contenido y la presenta como un frontier system completo. Lo más importante no es el nombre de un producto, sino los siguientes principios, aplicables a otros contextos:

\begin{enumerate}
  \item \textbf{Prioridad al resultado de extremo a extremo}: signal, platform decision, lawful report y victim protection son etapas distintas y requieren un handoff claro entre ellas;
  \item \textbf{Separación en capas de lo known y lo unknown}: hash matching procesa a alta velocidad los objetos ya conocidos y predictive detection amplía la cobertura de candidatos desconocidos; ambos requieren policy y revisión;
  \item \textbf{Diseño conjunto del modelo y la cola}: precision, recall, calibration, threshold, review capacity y moderator wellbeing no pueden validarse por separado;
  \item \textbf{La gobernanza de datos es una capacidad central}: trusted source, provenance, segregation, retention, label protocol y audit determinan si el modelo puede desplegarse;
  \item \textbf{Safety by Design abarca todo el ciclo de vida}: development, deployment, maintenance, red teaming y transparency tienen cada uno un owner definido;
  \item \textbf{La privacidad es un objetivo de arquitectura}: E2EE, metadata, endpoint y network signal tienen cada uno sus propios límites de derechos y de seguridad; deben minimizarse y admitir apelación;
  \item \textbf{Las señales más tempranas exigen acciones más mesuradas}: conversation y graph analysis aportan alertas tempranas, pero las acciones deben ser escalonadas, reversibles y conservar la human review;
  \item \textbf{Los roles del ecosistema no son intercambiables}: plataformas, proveedores de tecnología, organismos receptores de reportes, investigadores y servicios de apoyo forman un ciclo cerrado de protección mediante contract y audit.
\end{enumerate}

\begin{importantbox}{Tarea de seguridad de plataformas: diseñar un ciclo cerrado verificable}
Elige un producto que permita a los usuarios subir contenido o enviar mensajes privados y entrega un diseño de una página: dibuja la máquina de estados de detección, revisión, actuación y reporte; define los campos de known-match y de predictive signal; establece las restricciones de precision/recall y de reviewer-capacity; redacta las políticas de privacy, retention, appeal, incident y provider-outage. No se permite dibujar solo un recuadro de «AI moderation»: debe indicarse el owner y la evidencia de cada handoff.
\end{importantbox}

\subsection{Lecturas adicionales}

{\small
\begin{itemize}[nosep,leftmargin=*]
  \item \href{https://www.thorn.org/solutions/for-platforms/}{\textbf{Thorn platform solutions}}: punto de acceso oficial a Safer detection, Match y Predict.
  \item \href{https://info.thorn.org/hubfs/thorn-safety-by-design-for-generative-AI.pdf}{\textbf{Thorn Safety by Design for Generative AI}}: documento técnico sobre principios para el ciclo de vida de la IA generativa.
  \item \href{https://www.missingkids.org/gethelpnow/cybertipline/cybertiplinedata}{\textbf{NCMEC CyberTipline data}}: criterios de los informes anuales y descripción de tendencias.
  \item \href{https://www.missingkids.org/netsmartz/home}{\textbf{NCMEC NetSmartz}}: recursos de educación preventiva dirigidos a familias, escuelas y niños.
  \item \href{https://nvlpubs.nist.gov/nistpubs/ai/NIST.AI.600-1.pdf}{\textbf{NIST AI 600-1}}: Generative AI Profile.
  \item \href{https://nvlpubs.nist.gov/nistpubs/ai/NIST.AI.100-4.pdf}{\textbf{NIST AI 100-4}}: synthetic-content risk technical approaches.
\end{itemize}
}

\end{document}
